%!TEX TS-program = xelatex
%!TEX encoding = UTF-8 Unicode
% !Mode:: "TeX:UTF-8"
\documentclass[10pt,fancyhdr,UTF8, fontset = mac]{ctexbook}
\usepackage[centering,paperwidth=180mm,paperheight=260mm, body={390pt,600pt}]{geometry}
%\usepackage[centering,paperwidth=180mm,paperheight=230mm, body={390pt,530pt},showframe]{geometry}

%\documentclass[a4paper,11pt, svgnames]{scrbook} %如果不加oneside，则奇偶页的边距会不同，以方便装订
%\usepackage[a4paper, top=2cm, bottom=2cm, left=2cm, right=2cm]{geometry}
\usepackage{ulem}

\usepackage[scheme=plain]{ctex}
%\usepackage{xeCJK}

\usepackage{fontspec, xunicode, xltxtra, fancybox, setspace}
%\usepackage{ fontspec, xunicode, xltxtra, fancybox, setspace,xcolor}
\usepackage[dvipsnames,table,xcdraw, svgnames]{xcolor}
\usepackage{pifont} % special symbols like circle number etc. 

\usepackage{amsmath, amsthm, amsfonts, amssymb, mathrsfs} %数学公式相关 amssymb用于斜体字符，如varnothing,  masthm用于定理, mathrsfs用于花体字母
\newtheorem{theorem}{\hskip 2em \cjkem 定理}[chapter]
\newtheorem{lemma}{\hskip 2em \cjkem 引理}[chapter]
\newtheorem{definition}{\hskip 2em 定义}[chapter]
% \newtheorem{example}{\hskip 2em \cjkem 例} [chapter] % 该用tcolorbox实现
\renewcommand{\proofname}{\hskip 2em \cjkem 证明}
%\usepackage[sort]{natbib}

\usepackage[ruled, lined, linesnumbered, algochapter]{algorithm2e}
\renewcommand{\algorithmcfname}{算法}
\SetKwInput{KwIn}{输入}
\SetKwInput{KwOut}{输出}

\usepackage{colortbl, booktabs, multirow, longtable, makecell} %表格处理相关的Package
\usepackage{array}
\newcolumntype{x}[1]{>{\centering\let\newline\\\arraybackslash\hspace{0pt}}p{#1}}

\usepackage{tikz, flowchart} % tikz绘图
\usepackage{pgfplots}
\usetikzlibrary{decorations.pathmorphing} 
\usetikzlibrary{decorations.pathreplacing}
\usetikzlibrary{decorations.markings}
\usetikzlibrary{calc, arrows, arrows.meta, shapes, shadows,  shapes.geometric, positioning}
\usetikzlibrary{topaths}
\usetikzlibrary{matrix}
\usetikzlibrary{calc}
\usepackage{forest}

\usepackage[tikz]{bclogo}  %see http://mirrors.ctan.org/graphics/bclogo/README
\DeclareGraphicsRule{.mps}{eps}{*}{} %解决xelatex处理bclogo时的mps问题
\newcommand{\infobox}[2]{ 
    \begin{bclogo}[couleur=yellow!10, logo=\bcfleur, ombre=true]{#1}
    #2
    \end{bclogo}
}
\newcommand{\warnbox}[2]{ 
    \begin{bclogo}[couleur=yellow!10, logo=\bctakecare, ombre=true]{#1}
    #2
    \end{bclogo}
}

\newcommand{\shellbox}[1]{ 
    \begin{bclogo}[couleur=gray!10, logo=\bccrayon, ombre=true]{命令}
    #1
    \end{bclogo}
}

\usepackage{indentfirst}  %首行缩进
\setlength{\parindent}{2em} 

\usepackage{fancyhdr}
\thispagestyle{empty}

% 设置 plain style 的属性
\fancypagestyle{plain}{%
		\fancyhf{} % 清空当前设置

		% 设置页眉 (head)
		\fancyhead[RE]{\leftmark} % 在偶数页的右侧显示章名
		\fancyhead[LO]{\rightmark} % 在奇数页的左侧显示小节名
		\fancyhead[LE,RO]{~\thepage~} % 在偶数页的左侧，奇数页的右侧显示页码

		% 设置页脚：在每页的右下脚以斜体显示书名
		%\fancyfoot[RO,RE]{ {\it Written by Summer XIA}, Email: {\it xiat@ruc.edu.cn}}

		\renewcommand{\headrulewidth}{0.7pt} % 页眉与正文之间的水平线粗细
		\renewcommand{\footrulewidth}{0pt}
}

\pagestyle{fancy} % 选用 fancy style
\fancyhf{}
\fancyhead[RE]{\normalfont\small\rmfamily\nouppercase{\leftmark}}
\fancyhead[LO]{\normalfont\small\rmfamily\nouppercase{\rightmark}}
\fancyhead[LE,RO]{\thepage}
%\fancyfoot[LE,LO]{\small Book Title}


%%%%%%%%%% 一些重定义 %%%%%%%%%%
\usepackage{makeidx}
\usepackage{imakeidx}
\makeindex[title=索引, columns=2]
\renewcommand{\contentsname}{目录}     % 将Contents改为目录
\renewcommand{\indexname}{索引}
\renewcommand{\figurename}{图}
\renewcommand{\tablename}{表}
%\usepackage[figurename=图, tablename=表]{caption}
\renewcommand{\appendixname}{附录}
\renewcommand{\bibname}{参考文献}
\usepackage{caption} 
\captionsetup[table]{skip=6pt} % 设置表格标题和表格之间的间距大小

\usepackage{xurl} % 支持URL换行
\usepackage{diagbox} % 斜线表格

\let\oldbibliography\thebibliography
\renewcommand\thebibliography[1]{
	\oldbibliography{#1}
	\setlength{\parskip}{0pt}
	\setlength{\itemsep}{0pt plus 0.2ex}
}

\newcommand{\tab}{\phantom{o}\hspace{2ex}}


%\usepackage{draftwatermark}
%\SetWatermarkText{Xia Tian}%设置水印文字
%\SetWatermarkLightness{0.9}%设置水印亮度
%\SetWatermarkScale{0.5}%设置水印大小

\usepackage[bookmarks=true, bookmarksnumbered=true]{hyperref} %生成pdf索引
\usepackage{bookmark}
\hypersetup{colorlinks, citecolor=black,  filecolor=black, linkcolor=black,  urlcolor=black}

\usepackage{kpfonts}
\usepackage[bf, explicit]{titlesec} %对标题的式样进行设置
\newcommand*\chapterlabel{}
\titlespacing*{\chapter}{0pt}{50pt}{-60pt}
\renewcommand{\chaptername}{}
\titleformat{\section}[hang]{\color{SteelBlue}\LARGE\bfseries}{\color{SteelBlue}{\thesection}}{.8em}{#1}
\titleformat{\subsection}[hang]{\color{SteelBlue}\Large\bfseries}{\color{SteelBlue}{\thesubsection}}{.8em}{#1}
\titleformat{\subsubsection}[hang]{\color{SteelBlue}\large\bfseries}{\color{SteelBlue}{\thesubsubsection}}{.8em}{#1}

\titleformat{\paragraph}[hang]{\large\bfseries}{\theparagraph}{.8em}{\hspace{1em}$\blacksquare$ #1}

%\newcommand*{\hei}{\fontfamily{FZLanTingHeiS-H-GB}\selectfont}
%\DeclareTextFontCommand{\texthei}{\hei}

%\setCJKfamilyfont{SourceHanSansCN-Light.ttf}
%\setCJKmainfont{Source Han Sans SC VF Light} %方正字体，也可以改成：微软雅黑
%\setCJKsansfont{Source Han Sans SC VF Light}
%\setCJKfamilyfont{SourceHanSansCN-Light}
\setCJKmainfont{SourceHanSansCN-Light}


\newenvironment{block}[2]{\color{#1} \vspace{0.3cm} #2 \vspace{0.5cm} }{}


% Fix problem in tcolorbox:  ! LaTeX Error: Not in outer par mode.
% Add \usepackage{float} and use the option [H] in the figure.
\usepackage{float}
% \usepackage{multicol} % 解决tcolorbox跨页问题
\usepackage[many]{tcolorbox}
\tcbset{colback=white,colframe=black!60,fonttitle=\bfseries,
  coltitle=black}
\tcbuselibrary{theorems}
\tcbuselibrary{skins}
\newtcbtheorem[number within=chapter]{example}{例 }%
{ enhanced,
  toptitle=3mm,bottomtitle=5mm,
  breakable, pad at break=2mm,break at=-\baselineskip/0pt,
  colback=white,colbacklower=yellow,
  colbacktitle=black!5,colframe=black!60, fonttitle=\bfseries, arc=1mm,
  before skip=10pt,after skip=10pt,
  before upper={\parindent 2em}, before lower={\parindent 2em},
  arc=4mm,outer arc=1mm,
  attach boxed title to top left= {xshift=3.2mm,yshift=-0.10mm},
  boxed title style={skin=enhancedfirst jigsaw, arc=2mm,bottom=-1mm},
  segmentation style={draw,solid,decorate,decoration={coil,aspect=0,segment length=10.1mm}}
}{ex} % ex表示例子的label的开始字符

\usepackage{minted} %compile： lualatex/xelatex -shell-escape spark.tex
\definecolor{bg}{rgb}{0.975,0.975,0.975}
\setminted{encoding=utf-8} %注意字体，如果设置了CJKmonofont，会出现乱码
\setminted{linenos=true,tabsize=4, xleftmargin=1mm, breaklines, fontfamily=courier,fontsize=\small}
\setminted{bgcolor=bg, numbersep=2mm}
\usemintedstyle{pastie}

\usepackage{enumitem}
%\setlist[itemize]{align=parleft,left=0pt..1em, leftmargin=3em}
\setenumerate[1]{itemsep=0pt,partopsep=0pt,parsep=\parskip,topsep=0pt,}
\setitemize[1]{itemsep=0pt,partopsep=0pt,parsep=\parskip,topsep=0pt}
\setdescription{itemsep=0pt,partopsep=0pt,parsep=\parskip,topsep=0pt}

%自定义的一些命令，方便使用
\newcommand*\circled[1]{\tikz[baseline=(char.base)]{
  \node[shape=circle,draw,inner sep=1.5pt] (char) {#1};}}
%加入upcite，实现上标引用文献
\newcommand{\upcite}[1]{\textsuperscript{\textsuperscript{\cite{#1}}}} 

% 利用pifont在heading前加上方框符号
\newcommand{\heading}[1]{  {\vspace{4pt} \color{darkgray}\ding{111} \textbf{#1}} \vspace{2pt}}

% 重定义数学向量符号的形式
\renewcommand{\vec}[1]{\symbf{#1}}  
%\renewcommand{\vec}[1]{\overrightarrow{#1}} 

% 设置表格的行高
\renewcommand{\arraystretch}{1.5}
\renewcommand{\toprule}{\Xhline{0.8pt}}
\renewcommand{\midrule}{\Xhline{0.5pt}}
\renewcommand{\bottomrule}{\Xhline{0.7pt}}


\setstretch{1.35} % 利用setspace包提供的命令，设置行间距
\begin{document}
    \title{
    \huge{\textcolor{blue!80}{Python与数据结构}} \vspace{0.5cm} \\
    \large{\textcolor{blue!45}{Data Structure}} \vspace{1cm}
    \\ \vspace{2cm}
\begin{tikzpicture}[
node/.style = {circle, draw, fill=blue!20!black, inner sep=0pt, minimum size=8mm},
edge/.style = {thick}
]

% 定义节点
\node[node] (A) at (0,0) {};
\node[node] (B) at (2,1) {};
\node[node] (C) at (2,-1) {};
\node[node] (D) at (-2,1) {};
\node[node] (E) at (-2,-1) {};

% 定义边
\draw[edge] (A) -- (B);
\draw[edge] (A) -- (C);
\draw[edge] (A) -- (D);
\draw[edge] (A) -- (E);
\draw[edge] (B) -- (C);
\draw[edge] (D) -- (E);

\end{tikzpicture}
\vspace{4cm}
\\~ \\
}
\author{
	XiaTian \\
	Renmin University of China
}

\pagestyle{empty} % 关闭页码输出
\maketitle
\newpage
~ 
\pagestyle{empty} % 关闭页码输出
       % include title page

    \frontmatter %用于生成罗马计数的前言
    \pagestyle{plain}

	\titleformat{\chapter}
	  {\gdef\chapterlabel{}
	   \normalfont\sffamily\Huge\bfseries\scshape}
	  {\gdef\chapterlabel{\thechapter\ }}{0pt}
	  {
		\begin{tikzpicture}[remember picture,overlay]
		\node[yshift=-4cm] at (current page.north west)
		  {\begin{tikzpicture}[remember picture, overlay]
		  	\draw[fill=LightSkyBlue!20, draw=black!20] (0,0) rectangle
			  (\paperwidth,4cm);
			\draw node[anchor=east,
				  xshift=.9\paperwidth, 
			      yshift=0.2cm,
				  rectangle,
				  fill=white, 
				  draw=black!20,
				  rounded corners=10pt,
				  inner sep=11pt,
				  minimum width=6cm](name)
				  {  {\chapterlabel} \color{black}  #1   };
		   \end{tikzpicture}
		  };
	   \end{tikzpicture}
	   \vspace{2em}
	  }

    \include{preface}

    %generate book content list
    \setcounter{secnumdepth}{4}  % 编号定位到subsubsection
    \setcounter{tocdepth}{4}
    \small % 目录字体变小
    \tableofcontents
    \normalsize{}

    \mainmatter

	\titleformat{\chapter}
	  {\gdef\chapterlabel{}
	   \normalfont\sffamily\Huge\bfseries\scshape}
	  {\gdef\chapterlabel{\thechapter\ }}{0pt}
	  {
		\begin{tikzpicture}[remember picture,overlay]
		\node[yshift=-5cm] at (current page.north west)
		  {
		  	\begin{tikzpicture}[remember picture, overlay]
				\draw[fill=red!5, draw=black!20] (0,0) rectangle (\paperwidth,5cm);
				\draw node[scale=1.3, 
						rounded corners=5pt, 
						line width=1pt,  
						minimum width=2cm, 
						minimum height=2cm] at (3.5,2.5) { 
						 \color{black!80}~$\chapterlabel$ $\hookleftarrow$ 
					 	};
				\draw node[anchor=east,
							xshift=.9\paperwidth,
							yshift=0.2cm,
							rectangle,
					  		rounded corners=20pt,
					  		inner sep=11pt,draw=black!20,
					  		fill=white, 
					  		minimum width=7cm](name) {  \color{black!80}  #1   };
		   \end{tikzpicture}
		  };
	   \end{tikzpicture}
		\vspace{4em}
	  }

     %\include{chapter/introduction}
     \chapter{线性结构}

	 \chapter{栈（Stack）：后进先出的智慧}

在学习了通用的“线性表”之后，我们进入一类特殊的线性表——\emph{受限的线性表}。所谓受限，并不是功能变弱了，而是通过规则的约束，让它在特定场景下变得极其强大。栈就是这样一种结构。

\section{什么是栈？——后进先出的世界}

栈（Stack）是一个特殊的线性表，它只允许在\emph{一端}进行插入和删除操作。这一端称为\emph{栈顶}，另一端称为\emph{栈底}。不含任何元素的栈称为\emph{空栈}。

\subsection{核心原则：后进先出（LIFO）}

栈的精髓是 \emph{LIFO（Last In First Out）}，即后进先出。最后放入栈中的元素，最先被取出。

\textbf{生活中的类比}：

\begin{itemize}
    \item \emph{叠盘子}：你洗好一个盘子就叠在最上面；要用盘子时也是从最上面拿。最后放上去的盘子最先被拿走。
    \item \emph{步枪弹夹}：子弹一颗一颗压进去（入栈），射击时最上面那颗最先出膛（出栈）。
    \item \emph{浏览器后退按钮}：你访问的页面按顺序被压入一个栈，点击“后退”时，最后访问的页面最先被退出。
\end{itemize}

\subsection{为什么是“受限”的？}

与普通线性表相比，栈的操作被严格限制在一端。这种限制带来了两个好处：

\begin{itemize}
    \item \emph{逻辑清晰}：我们只关心栈顶元素，无需考虑中间元素的随机访问。
    \item \emph{实现简单}：只需要维护一个栈顶指针，插入和删除操作都集中在同一端。
\end{itemize}

\section{栈的基本操作}

栈的操作非常简洁，主要就是两个动作：

\begin{itemize}
    \item \emph{入栈（Push）}：将一个元素放到栈顶，成为新的栈顶。
    \item \emph{出栈（Pop）}：取出当前栈顶元素，栈顶指针下移，原来的次顶元素成为新栈顶。
\end{itemize}

\subsection{Python 中的栈实现}

Python 内置的 \texttt{list} 类型就是一个天然的栈：\texttt{append()} 对应入栈，\texttt{pop()} 对应出栈。我们可以用它来模拟浏览器的前进后退：

\begin{minted}{python}
# 初始化一个空栈，记录浏览历史
history_stack = []

# 模拟访问页面（入栈）
history_stack.append("首页")
history_stack.append("搜索页")
history_stack.append("详情页")
print(f"当前浏览路径: {history_stack}")

# 模拟点击“返回”按钮（出栈）
if history_stack:
    last_page = history_stack.pop()
    print(f"离开: {last_page}")
    print(f"回到上一页: {history_stack[-1] if history_stack else '空'}")
\end{minted}

运行这段代码，你会看到“详情页”被弹出，当前页面回到了“搜索页”。这正是后进先出的体现。

\section{栈的两种存储结构}

根据底层存储方式的不同，栈可以分为顺序栈和链栈。

\subsection{顺序栈}

顺序栈使用一组连续的存储单元（如数组）来存放元素，栈顶通常用一个整型变量 \texttt{top} 来指示。

\begin{itemize}
    \item \emph{优点}：访问速度快，实现简单。
    \item \emph{缺点}：需要预先分配空间，当栈满时无法再插入新元素，会发生\emph{栈溢出（Stack Overflow）}。
\end{itemize}

\begin{quote}
\emph{小知识}：“栈溢出”这个词大家可能经常在编程错误中见到，它就是因为递归调用太深或局部数据太大，导致系统分配的栈空间不够用了。
\end{quote}

\subsection{链栈}

链栈使用链表结构，每个节点包含数据和指向下一个节点的指针。栈顶就是链表的头节点。

\begin{itemize}
    \item \emph{优点}：动态分配内存，理论上不会满（除非内存耗尽）。
    \item \emph{缺点}：每个节点需要额外的指针空间，访问速度略慢于顺序栈。
\end{itemize}

现代高级语言（如 Java、Python）的虚拟机在管理方法调用时，底层大量使用了链式栈的思想，以便灵活应对复杂的调用层次。

\section{栈的经典应用}

栈不仅仅是一种存储数据的方式，它代表了一种\emph{“保存现场、稍后回来”}的思维模式。下面我们看看栈在现实和算法中的几个典型应用。

\subsection{撤销操作（Undo）}

在文字处理软件（如 Word）中，你的每一次编辑（输入文字、删除、格式化）都会被压入一个“撤销栈”。当你按下 Ctrl+Z 时，系统从栈中弹出最近的操作，并执行反向动作，从而恢复到之前的状态。这正是后进先出思想的体现——最后做的操作最先被撤销。

\subsection{括号匹配}

编译器如何判断代码中的括号是否成对匹配？比如 \texttt{(( ))} 是合法的，而 \texttt{({ )}} 是不合法的。算法很简单：

\begin{itemize}
    \item 从左到右扫描字符串，遇到左括号（\texttt{(}、\texttt{\{}、\texttt{[}）就将其\emph{入栈}。
    \item 遇到右括号（\texttt{)}、\texttt{\}}、\texttt{]}）时，从栈顶\emph{弹出一个左括号}，检查是否匹配。
    \item 如果栈为空却遇到右括号，或者最后栈中还有剩余左括号，说明匹配失败。
\end{itemize}

这个算法只需要一次遍历，时间复杂度 O(n)，空间复杂度 O(n)（最坏情况全是左括号）。很多编程面试题都喜欢考察这个题目。

\subsection{表达式求值（逆波兰表达式）}

计算器在处理 \texttt{1 + 2 * 3} 时，并不是简单地从左到右计算，因为乘法的优先级高于加法。栈可以帮助我们暂存操作数和运算符，从而正确处理优先级。

一种经典方法是先将中缀表达式转换为后缀表达式（逆波兰表达式），然后用一个栈来求值。例如，\texttt{1 + 2 * 3} 对应的后缀表达式是 \texttt{1 2 3 * +}，求值过程如下：

\begin{itemize}
    \item 遇到操作数 \texttt{1}、\texttt{2}、\texttt{3} 依次入栈。
    \item 遇到运算符 \texttt{*}，弹出两个操作数 \texttt{2} 和 \texttt{3}，计算 \texttt{2 * 3 = 6}，将 \texttt{6} 入栈。
    \item 遇到运算符 \texttt{+}，弹出 \texttt{6} 和 \texttt{1}，计算 \texttt{1 + 6 = 7}，得到结果。
\end{itemize}

整个过程只需要一个栈，非常简洁。

\subsection{函数调用与递归}

在程序运行时，每当调用一个函数，系统就会在内存的“调用栈”上为该函数分配一块区域（称为栈帧），用来存放局部变量、参数和返回地址。当函数执行完毕，对应的栈帧被弹出，控制权交还给调用者。

递归函数正是利用了栈的这种特性：每次递归调用都会压入一个新的栈帧，返回时再弹出。如果递归深度过大，就会导致栈溢出。

\section{深度探索：迷宫寻路与回溯}

栈在人工智能和路径规划中的一个经典应用是\emph{迷宫寻路}。假设我们在一个二维迷宫中，从起点出发，要找到一条通往终点的路径。我们可以用栈来记录走过的路径，并在遇到死胡同时回溯。

\textbf{基本思路}：

\begin{enumerate}
    \item 从起点出发，将当前坐标压入栈。
    \item 探测四个方向（上、下、左、右）：
    \begin{itemize}
        \item 如果某个方向可以走（未走过且不是墙），就将新坐标压栈，并标记为已走。
        \item 如果所有方向都不可走，说明当前是死胡同，则\emph{出栈}（回溯）到上一个位置，尝试其他方向。
    \end{itemize}
    \item 重复步骤 2，直到到达终点或栈空（无路可走）。
\end{enumerate}

这种“试探-回溯”的过程，正是\emph{深度优先搜索（DFS）}的核心思想，而栈则是实现这一思想的天然工具。

\begin{quote}
为什么栈适合回溯？\\
因为回溯需要回到最近访问的点，而栈的后进先出特性恰好能记住最近的选择。每一步的决策都被压入栈，当发现此路不通时，只需弹出栈顶，就能回到上一个分叉点。
\end{quote}

\section{栈与队列的对比}

在数据结构中，栈和队列经常被放在一起讨论，因为它们都是受限的线性表，但规则恰好相反：

\begin{table}[h]
\centering
\begin{tabular}{|l|l|l|}
\hline
\textbf{特性} & \textbf{栈（Stack）} & \textbf{队列（Queue）} \\
\hline
原则 & 后进先出（LIFO） & 先进先出（FIFO） \\
操作端 & 同一端（栈顶） & 两端（队尾插入，队头删除） \\
适用场景 & 回溯、撤销、递归 & 排队、任务调度、广度优先搜索 \\
\hline
\end{tabular}
\caption{栈与队列的对比}
\end{table}

简单来说，栈是为了“记住过去”（回溯），而队列是为了“保证公平”（先进先出）。

\section{自测练习}

为了巩固对栈的理解，这里给出几道练习题。你可以尝试用笔算，也可以动手写代码验证。

\begin{enumerate}
    \item \textbf{出栈序列判断} \\
    若入栈序列为 \texttt{A, B, C}，且允许在任意时刻出栈，下列哪个序列是不可能出现的出栈序列？ \\
    A. \texttt{C, B, A} \\
    B. \texttt{B, A, C} \\
    C. \texttt{C, A, B} \\
    D. \texttt{A, B, C}

    \item \textbf{括号匹配加强版} \\
    给定一个只包含 \texttt{(}、\texttt{)}、\texttt{\{}、\texttt{\}}、\texttt{[}、\texttt{]} 的字符串，判断它是否有效。要求使用栈实现。

    \item \textbf{逆波兰表达式求值} \\
    输入一个逆波兰表达式（后缀表达式），例如 \texttt{["2","1","+","3","*"]}，计算其结果。

    \item \textbf{模拟浏览器的前进后退} \\
    设计一个浏览器历史记录类，支持访问新页面、后退、前进操作。用两个栈来实现。

    \item \textbf{最小栈} \\
    设计一个栈，除了支持常规的 push、pop 操作外，还要能在常数时间内获取栈中的最小元素。

    \item \textbf{用栈实现队列} \\
    使用两个栈来实现一个队列，支持入队和出队操作。

    \item \textbf{回文链表检测} \\
    给定一个单链表，判断它是否回文（例如 \texttt{1->2->2->1}）。可以借助栈实现：先遍历一遍将值入栈，再第二次遍历时与栈顶比较。

    \item \textbf{迷宫问题} \\
    用栈实现一个简单的迷宫求解程序，输出从起点到终点的路径。
\end{enumerate}

\section*{本章小结}

栈是一种“简约而不简单”的数据结构。它在计算机底层（函数调用）、算法（深度优先搜索）和日常应用（撤销、括号匹配）中无处不在。掌握了栈，你就掌握了程序如何“记住过去”的方法。

在下一章中，我们将学习另一种受限的线性表——\emph{队列}，看看它又是如何实现“公平”的。



    \appendix
    %\include{history}       % include first appendix
    %\include{tables}        % include second appendix, etc.
    
    \backmatter

    \titleformat{\chapter}
	  {\gdef\chapterlabel{}
	   \normalfont\sffamily\Huge\bfseries\scshape}
	  {\gdef\chapterlabel{\thechapter\ }}{0pt}
	  {
		\begin{tikzpicture}[remember picture,overlay]
		\node[yshift=-4cm] at (current page.north west)
		  {\begin{tikzpicture}[remember picture, overlay]
		  	\draw[fill=blue!10, draw=black!20] (0,0) rectangle
			  (\paperwidth,4cm);
			\draw node[anchor=east,xshift=.9\paperwidth, 
			      yshift=0.2cm,
				  rectangle,
				  fill=white, draw=black!20,
				  rounded corners=10pt,inner sep=11pt,
				  minimum width=6cm](name)
				  {  {\chapterlabel} \color{black}  #1   };
		   \end{tikzpicture}
		  };
	   \end{tikzpicture}
	   \vspace{2em}
	  }

    \renewcommand\bibname{参考文献}
	\bibliographystyle{chinesebst}
    \bibliography{refs}
    \printindex
    %\include{chapter/thanks}
\end{document}

